\section{团队协作与持续演进}

\subsection{Roles \& responsibilities 矩阵}

成功的模仿学习系统需要 alignment、ops、data、legal 四个团队协作。下面的矩阵列出了每个角色的 output 和 checkpoint：

\begin{table}[H]
\centering
\begin{tabular}{p{0.22\textwidth} p{0.28\textwidth} p{0.28\textwidth} p{0.1\textwidth}}
\toprule
角色 & 交付物 & 核心 checkpoint & 频率 \\
\midrule
Alignment Engineer & Evidence matrix, policy checkpoint & KL drift < θ, preference score & per model release \\
Ops/SRE & Drift dashboard, latency log & alert resolution time & daily \\
Data Engineer & Demo catalog, metadata & coverage report & weekly \\
Legal/Compliance & Guardrail doc, review minutes & high-risk prompt gating & per sprint \\
\bottomrule
\end{tabular}
\caption{模仿学习团队的责任矩阵}
\end{table}

\begin{warningbox}{不要把责任模糊化}
如果某条 guardrail 既不是 ops 的责任也不是 compliance 的责任，就会在 incident 中无人响应。要用矩阵明确 owner，并把流程写入 team handbook。
\end{warningbox}

\subsection{持续演进与研究方向}

模仿学习在 research 轨道上可以迈向 sim-to-real、multi-agent collaboration、自动化 reward inference。Slide 02-11 强调了 multi-agent scenario 中的 cooperative imitation、competitive imitation 以及 dynamic gating。

\begin{importantbox}{Future experiments 的关键}
每次 research experiment 都要陪同 evaluation matrix，确认实验结束后 metric 是否回到 baseline；如果 metrics drift 严重，就要延迟 production rollout。
\end{importantbox}

\subsection{本章小结}

团队协作与 research planning 构成持续演进的双车轴：责任矩阵保证每次 release 有人在看数据，研究方向则通过 structured experiment 让模仿学习 pipelines 在多 agent、sim-to-real 的未来场景中稳健。

